% Exercise 2.3-6

Since the loop invariant of \textsc{Insertion-Sort} guarantees that we have
$A[1 \ltwodots j - 1]$ sorted, we can indeed use the \textsc{Binary-Search}
algorithm of 2.3-5 to locate the position at which $A[j]$ belongs, reducing the
number of \emph{comparisons} for that insertion from $\Theta(j)$ to
$\Theta(\lg{j})$.

This does not, however, improve the worst-case running time to
$\Theta(n\lg{n})$. The \textbf{while} loop of lines 5-7 does not only compare
elements; it also shifts each element of $A[1 \ltwodots j-1]$ that exceeds
$A[j]$ one position to the right. Binary search locates the insertion point but
performs none of that movement, and in the worst case, when $A$ is in
decreasing order, all $j - 1$ elements must still be moved. Summing over every
iteration of the outer loop gives
\[
  \sum_{j = 2}^n \Theta(j) = \Theta(n^2)
\]
element movements regardless of how the insertion point is found, so the
worst-case running time of \textsc{Insertion-Sort} remains $\Theta(n^2)$.
